\documentclass[11pt,a4paper]{article}
\usepackage[T1]{fontenc}
\usepackage[utf8]{inputenc}
\usepackage{lmodern,amsmath,amssymb,amsthm,booktabs}
\usepackage[margin=27mm]{geometry}
\usepackage[hidelinks]{hyperref}
\hypersetup{pdftitle={Time Refinement and an Action Scale},pdfauthor={navstokgap research working notes}}
\setlength{\emergencystretch}{3em}
\newtheorem{proposition}{Proposition}
\title{Time Refinement and an Action Scale}
\author{navstokgap research working notes}
\date{7 September 2026}
\begin{document}
\maketitle
\begin{abstract}
Free-particle kernels give an exact test of refinement-based proposals for an
action scale. Integrating over an inserted position preserves a Gaussian family;
time homogeneity makes its variance linear in duration. The proportionality
constant remains a model parameter. Fixed-endpoint bridges concentrate on the
Newtonian straight path as that parameter vanishes. We use these standard
calculations to isolate the physical selection question in the proposed
connection between geometric subdivision, renormalisation and quantum mechanics.
\end{abstract}

\section{A source-driven experiment}
The experiment is to refine a time interval, integrate out the inserted
positions, and track the remaining parameters. Rivero's cone discussion
\cite{RiveroCone1999} motivates retaining rescaled data as sections approach;
his path-integral proposal \cite{RiveroFeynman1998} motivates a finite blocking
map. Brouder's rooted-tree correspondence \cite{Brouder1999} supplies an
algebraic route to nonlinear composition, and the tangent-groupoid construction
\cite{CarinenaEtAl1999} supplies a separate quantum-completion reading.

This note works on the free line with mass $m>0$. Positions have units of
length, $s,t,T>0$ are durations, and $\kappa$ has units of action.
An identification $\kappa=\hbar=h/(2\pi)$ is a physical interpretation of a
selected oscillatory model. The present calculations classify a free family;
the selection of its physical member is the open reconstruction question.

\section{Eliminating an inserted position}
The normalized positive kernel is
\begin{equation}\label{eq:kernel}
 G_t^\kappa(x)=\sqrt{\frac{m}{2\pi\kappa t}}
 \exp\left(-\frac{mx^2}{2\kappa t}\right),\qquad \kappa>0.
\end{equation}
It is a probability density on $\mathbb R$ with variance $\kappa t/m$;
the standard Brownian normalization is given in
\cite[\S4.1, (23)--(25)]{PitmanYor2018}.
Its exponent is minus the free straight-segment action divided by $\kappa$.
Here positivity defines a probability model; the oscillatory model enters
in Section~\ref{sec:quantum}.

\begin{proposition}[Exact blocking]\label{prop:block}
For $s,t,\kappa>0$ and $x,z\in\mathbb R$,
\[
 \int_{\mathbb R}G_t^\kappa(z-y)G_s^\kappa(y-x)\,dy
 =G_{s+t}^\kappa(z-x).
\]
\end{proposition}
\begin{proof}
Set $\mu=(tx+sz)/(s+t)$. Completing the square gives
\[
 \frac{(y-x)^2}{s}+\frac{(z-y)^2}{t}
 =\frac{(z-x)^2}{s+t}+\frac{s+t}{st}(y-\mu)^2.
\]
The remaining Gaussian integral is
$\sqrt{2\pi\kappa st/[m(s+t)]}$. Multiplication by the two prefactors in
\eqref{eq:kernel} gives exactly the prefactor at duration $s+t$.
\end{proof}

Consequently, for any finite partition $0=t_0<\cdots<t_n=T$,
integrating the product of the $n$ kernels over its $n-1$ internal positions
gives $G_T^\kappa(z-x)$. The action parameter is unchanged at every finite
refinement; this identity requires no continuum path measure.

\begin{proposition}[The Gaussian parameter left by composition]\label{prop:classify}
Let $(\mu_t)_{t\ge0}$ be a weakly continuous convolution semigroup of centered
Gaussian probability measures on $\mathbb R$, allowing a degenerate Gaussian,
with $\mu_0=\delta_0$. Then $\mu_t$ has variance $w(t)=at$ for one $a\ge0$.
Conversely every $a\ge0$ defines such a family.
\end{proposition}
\begin{proof}
Characteristic functions are $\widehat\mu_t(k)=\exp[-w(t)k^2/2]$.
Composition implies $w(s+t)=w(s)+w(t)$. Weak continuity, evaluated at $k=1$,
makes $w$ continuous. Additivity first yields $w(r)=r w(1)$ for nonnegative
rational $r$, and continuity extends this to every real $t\ge0$.
Nonnegative variance gives $a=w(1)\ge0$. These characteristic functions
verify the converse.
\end{proof}

At fixed mass set $\kappa=ma$; since $[a]=\text{length}^2/\text{time}$,
$[\kappa]=\text{action}$. All $\kappa>0$ satisfy the same composition law,
and $a=0$ supplies the identity convolution family. Positive-width densities
select $a>0$ as an additional premise. Even that subclass has
$\inf\{\kappa:\kappa>0\}=0$: positivity of each member and a uniform lower
bound over a class are distinct questions.

\section{A bridge to the Newtonian path}
Conditioning on endpoints gives an exact trajectory comparison. The standard
bridge construction is given in \cite[\S3.4]{PitmanYor2018}.
For $0<s<T$, define the normalized intermediate-position density
\[
 \rho_s^\kappa(y\mid x,z)=
 \frac{G_{T-s}^\kappa(z-y)G_s^\kappa(y-x)}{G_T^\kappa(z-x)}.
\]
Proposition~\ref{prop:block} makes its integral one. The same square completion
gives mean and variance
\begin{equation}\label{eq:bridge}
 q_{\rm cl}(s)=x+\frac{s}{T}(z-x),\qquad
 \sigma_s^2=\frac{\kappa s(T-s)}{mT}.
\end{equation}
The mean is the unique free Newtonian trajectory with these endpoints.
For any displacement tolerance $r>0$, Chebyshev's inequality gives
\[
 \Pr\{|Y_s-q_{\rm cl}(s)|\ge r\}
 \le\frac{\kappa s(T-s)}{mTr^2}\longrightarrow0.
\]
For any fixed finite set of internal times, the product-kernel bridge has these
marginals, and a union bound proves joint concentration on the sampled straight
path. This is a finite-dimensional classical-path limit in a positive-measure
model. Quantum phase-space correspondence is a further limiting problem.

The zero-variance convolution family describes a stationary position. General
free classical motion instead has the phase-space transition measure
\[
 P_t((q,p),dq'\,dp')=
 \delta_{q+pt/m}(dq')\,\delta_p(dp'),
\]
whose composition follows by substitution. This keeps the deterministic
comparison on its appropriate state space.

\section{The oscillatory family and its coefficient}\label{sec:quantum}
For each $\kappa>0$, the free oscillatory kernel is
\[
 K_t^\kappa(x)=\sqrt{\frac{m}{2\pi i\kappa t}}
 \exp\left(\frac{imx^2}{2\kappa t}\right),\qquad t>0.
\]
Use the unitary transform
$\mathcal Ff(k)=(2\pi)^{-1/2}\int_{\mathbb R}e^{-ikx}f(x)\,dx$.
The kernel's precise meaning is the tempered-distribution kernel of the Fourier
multiplier $\exp[-i\kappa t k^2/(2m)]$, with the square-root branch fixed by
continuation from positive real Gaussian variance. On $L^2(\mathbb R)$,
these multipliers define a strongly continuous unitary group for real $t$:
the group law follows by multiplication, unitarity by modulus one, and strong
continuity by dominated convergence. Its Hamiltonian in the convention
$i\kappa\partial_t\psi=H_\kappa\psi$ is
\[
 H_\kappa=-\frac{\kappa^2}{2m}\frac{d^2}{dx^2},
 \qquad D(H_\kappa)=H^2(\mathbb R).
\]
These are the usual free-particle formulas \cite[\S7.3]{Teschl2009}.
Operator composition defines the oscillatory convolution, avoiding an
unqualified interchange of conditionally convergent integrals.

Writing the positive family as $\exp[t\kappa\partial_x^2/(2m)]$ and continuing
$t$ to $it$ yields this unitary family. That continuation is a specified model
choice. Both constructions retain a free coefficient with action units.

\section{The selection problem now exposed}
The first task is to find a physical principle that selects a positive action
scale and explains its universality under composition of interacting systems.
Three parameters stay separate: temporal mesh size, the action scale $\kappa$,
and any finite-speed parameter $c^{-1}$. The present free nonrelativistic model
allows arbitrarily fine temporal meshes at each fixed $\kappa$.

The next source-guided experiment audits the two regulators in
\cite{RiveroFeynman1998}, starting with a finite-dimensional quadratic action
and its exact normalization. A second task extracts the hypotheses and limiting
topology of \cite{CarinenaEtAl1999}. Together they test whether refinement
consistency constructs a quantum completion, and which additional premise
selects its physical action parameter. Interactions and finite propagation
then become explicit changes of model.

\clearpage
\bibliographystyle{plain}
\bibliography{references/library}
\end{document}
